% GENERATED FILE. Edit canonical lessons, then run npm run generate:books.
% canonical-source-hash: fnv1a64:87627d49ec383d09
% canonical-lessons: RU-C509-vorota, RU-C509-vratar, RU-C509-napadayushchiy, RU-C509-gol, RU-C509-zabivat

\chapter{Goals and Goalkeepers}
\label{ch:ru-goals-and-goalkeepers}
\hlchaptermodality{509}

\begin{chapteropening}
\textbf{By the end of this chapter:} \emph{I can talk about scoring and saving goals.}

Complete the last lesson of chapter 509: I can talk about scoring and saving goals.
\end{chapteropening}

\section[\texorpdfstring{\ru{ворота} (voróta)}{ворота (voróta)}]{\ru{ворота} (voróta) — gates; the goal (in football)}

\label{lesson:RU-C509-vorota}



\begin{quote}
Before the new one: say the Russian for to cheer on, then the Russian for to rejoice.
\end{quote}

\subsection*{You'll want to know: \ru{ворота}}
\textbf{\ru{ворота}} (\emph{voróta}) — \textquotedblleft{}gates; the goal (in football)\textquotedblright{}.

\begin{quote}
\ru{Мяч} \ru{пролетел} \ru{мимо} \ru{ворот}.
\end{quote}

\begin{quote}
\emph{Myach proletél mímo vorót.}
\end{quote}

\begin{quote}
The ball flew past the goal.
\end{quote}

Plural only; genitive \ru{ворот}.

\begin{cousinweb}
\emph{\ru{ворота}} is inherited from Proto-Slavic \emph{vorta} (EtymDB), with the native East Slavic \emph{-\ru{оро}-}.
\end{cousinweb}

\begin{grammarlens}[title={What you've built}]
One more word for telling, explaining and describing at B1.
\end{grammarlens}

\subsection*{Your turn}
\begin{itemize}
\raggedright
  \item \emph{Say it:} \emph{voróta}
  \item \emph{Say it:} \emph{voróta}, once more
  \item \emph{Say it:} a sentence of your own with \emph{voróta}
  \item \emph{Recall:} say the Russian for to cheer on, then the Russian for to rejoice, then say \emph{voróta} again
\end{itemize}

\begin{tcolorbox}[breakable,colback=teal!4,colframe=teal!35!black,title={Before you move on}]
What does \ru{ворота} mean? (\textquotedblleft{}Gates; the goal (in football)\textquotedblright{}.) Say it once more.
\end{tcolorbox}
\section[\texorpdfstring{\ru{вратарь} (vratár)}{вратарь (vratár)}]{\ru{вратарь} (vratár) — a goalkeeper}

\label{lesson:RU-C509-vratar}



\begin{quote}
Before the new one: say the Russian for to rejoice, then the Russian for gates.
\end{quote}

\subsection*{You'll want to know: \ru{вратарь}}
\textbf{\ru{вратарь}} (\emph{vratár}) — \textquotedblleft{}a goalkeeper\textquotedblright{}.

\begin{quote}
\ru{Вратарь} \ru{поймал} \ru{мяч}.
\end{quote}

\begin{quote}
\emph{Vratár poymál myach.}
\end{quote}

\begin{quote}
The goalkeeper caught the ball.
\end{quote}

Masculine; genitive \ru{вратаря}.

\begin{cousinweb}
\emph{\ru{вратарь}} is built on \emph{\ru{врата}} with the ending \emph{-\ru{арь}} (EtymDB); \emph{\ru{врата}} was borrowed from Church Slavonic, the \emph{-\ru{ра}-} shape of Proto-Slavic \emph{vorta}, which gave native \emph{\ru{ворота}} 'gates' (EtymDB).
\end{cousinweb}

\begin{grammarlens}[title={What you've built}]
One more word for telling, explaining and describing at B1.
\end{grammarlens}

\subsection*{Your turn}
\begin{itemize}
\raggedright
  \item \emph{Say it:} \emph{vratár}
  \item \emph{Say it:} \emph{vratár}, once more
  \item \emph{Say it:} a sentence of your own with \emph{vratár}
  \item \emph{Recall:} say the Russian for to rejoice, then the Russian for gates, then say \emph{vratár} again
\end{itemize}

\begin{tcolorbox}[breakable,colback=teal!4,colframe=teal!35!black,title={Before you move on}]
What does \ru{вратарь} mean? (\textquotedblleft{}A goalkeeper\textquotedblright{}.) Say it once more.
\end{tcolorbox}
\section[\texorpdfstring{\ru{нападающий} (napadáyushchiy)}{нападающий (napadáyushchiy)}]{\ru{нападающий} (napadáyushchiy) — a forward, a striker}

\label{lesson:RU-C509-napadayushchiy}



\begin{quote}
Before the new one: say the Russian for gates, then the Russian for a goalkeeper.
\end{quote}

\subsection*{You'll want to know: \ru{нападающий}}
\textbf{\ru{нападающий}} (\emph{napadáyushchiy}) — \textquotedblleft{}a forward, a striker\textquotedblright{}.

\begin{quote}
\ru{Наш} \ru{нападающий} — \ru{самый} \ru{быстрый} \ru{в} \ru{команде}.
\end{quote}

\begin{quote}
\emph{Nash napadáyushchiy — sámyy býstryy v kománde.}
\end{quote}

\begin{quote}
Our striker is the fastest in the team.
\end{quote}

Declined like an adjective. The present participle of \emph{\ru{нападать}} — to attack.

\begin{cousinweb}
\emph{\ru{нападающий}} is the present participle of \emph{\ru{нападать}} 'to attack', built on \emph{\ru{на}-} and \emph{\ru{падать}} 'to fall'; EtymDB has no entry.
\end{cousinweb}

\begin{grammarlens}[title={What you've built}]
One more word for telling, explaining and describing at B1.
\end{grammarlens}

\subsection*{Your turn}
\begin{itemize}
\raggedright
  \item \emph{Say it:} \emph{napadáyushchiy}
  \item \emph{Say it:} \emph{napadáyushchiy}, once more
  \item \emph{Say it:} a sentence of your own with \emph{napadáyushchiy}
  \item \emph{Recall:} say the Russian for gates, then the Russian for a goalkeeper, then say \emph{napadáyushchiy} again
\end{itemize}

\begin{tcolorbox}[breakable,colback=teal!4,colframe=teal!35!black,title={Before you move on}]
What does \ru{нападающий} mean? (\textquotedblleft{}A forward, a striker\textquotedblright{}.) Say it once more.
\end{tcolorbox}
\section[\texorpdfstring{\ru{гол} (gol)}{гол (gol)}]{\ru{гол} (gol) — a goal}

\label{lesson:RU-C509-gol}



\begin{quote}
Before the new one: say the Russian for a goalkeeper, then the Russian for a forward.
\end{quote}

\subsection*{You'll want to know: \ru{гол}}
\textbf{\ru{гол}} (\emph{gol}) — \textquotedblleft{}a goal\textquotedblright{}.

\begin{quote}
\ru{В} \ru{первом} \ru{тайме} \ru{не} \ru{было} \ru{ни} \ru{одного} \ru{гола}.
\end{quote}

\begin{quote}
\emph{V pérvom táyme né bylo ni odnogó góla.}
\end{quote}

\begin{quote}
There wasn't a single goal in the first half.
\end{quote}

Masculine. \emph{\ru{Забить} \ru{гол}} — to score a goal.

\begin{cousinweb}
\emph{\ru{гол}} was borrowed from English \emph{goal} (EtymDB).
\end{cousinweb}

\begin{grammarlens}[title={What you've built}]
One more word for telling, explaining and describing at B1.
\end{grammarlens}

\subsection*{Your turn}
\begin{itemize}
\raggedright
  \item \emph{Say it:} \emph{gol}
  \item \emph{Say it:} \emph{gol}, once more
  \item \emph{Say it:} a sentence of your own with \emph{gol}
  \item \emph{Recall:} say the Russian for a goalkeeper, then the Russian for a forward, then say \emph{gol} again
\end{itemize}

\begin{tcolorbox}[breakable,colback=teal!4,colframe=teal!35!black,title={Before you move on}]
What does \ru{гол} mean? (\textquotedblleft{}A goal\textquotedblright{}.) Say it once more.
\end{tcolorbox}
\section[\texorpdfstring{\ru{забивать} (zabivát)}{забивать (zabivát)}]{\ru{забивать} (zabivát) — to score (a goal); to hammer in}

\label{lesson:RU-C509-zabivat}



\begin{quote}
Before the new one: say the Russian for a forward, then the Russian for a goal.
\end{quote}

\subsection*{You'll want to know: \ru{забивать}}
\textbf{\ru{забивать}} (\emph{zabivát}) — \textquotedblleft{}to score (a goal); to hammer in\textquotedblright{}.

\begin{quote}
\ru{Он} \ru{забивает} \ru{почти} \ru{в} \ru{каждом} \ru{матче}.
\end{quote}

\begin{quote}
\emph{On zabiváyet pochtí v kázhdom mátche.}
\end{quote}

\begin{quote}
He scores in almost every match.
\end{quote}

Imperfective; perfective \ru{забить} (\ru{забью}). Also to hammer in a nail.

\begin{cousinweb}
\emph{\ru{забивать}} is built on \emph{\ru{за}-} and the root of \emph{\ru{бить}} 'to strike'; EtymDB lists the verb without a source.
\end{cousinweb}

\begin{grammarlens}[title={What you've built}]
One more word for telling, explaining and describing at B1.
\end{grammarlens}

\subsection*{Your turn}
\begin{itemize}
\raggedright
  \item \emph{Say it:} \emph{zabivát}
  \item \emph{Say it:} \emph{zabivát}, once more
  \item \emph{Say it:} a sentence of your own with \emph{zabivát}
  \item \emph{Recall:} say the Russian for a forward, then the Russian for a goal, then say \emph{zabivát} again
\end{itemize}

\begin{tcolorbox}[breakable,colback=teal!4,colframe=teal!35!black,title={Before you move on}]
What does \ru{забивать} mean? (\textquotedblleft{}To score (a goal); to hammer in\textquotedblright{}.) Say it once more.
\end{tcolorbox}
